\section{Primero preguntar a las personas, después definir la tarea}

Si el objeto de estudio está directamente ligado a la vida comunitaria, la educación y la continuidad de la lengua, el equipo del modelo no puede imaginar las necesidades basándose solo en un benchmark. NLLB primero confirmó el problema mediante entrevistas y luego hizo participar a hablantes nativos en la traducción y la evaluación; este es el giro metodológico más importante de la clase, y también el que con más facilidad queda oculto tras los detalles del modelo.

\subsection{Cómo saber que se está resolviendo un problema real}

Esta sección pasa de los objetivos de cobertura a la definición del problema. Antes de mostrar cualquier pipeline de datos, Fan plantea un criterio de investigación: lo más importante es confirmar que el equipo está resolviendo un problema real, sobre todo cuando la tecnología está muy cerca de las personas. Esta pregunta no es un lema ético, sino el primer paso que determina los estándares de datos, las condiciones de éxito y el costo de los errores en lo que sigue.

\lecturefigure{slide-06-real-problem.jpg}{Confirmar, antes de modelar, que el problema existe para los usuarios reales}{Video de la clase de Stanford 00:05:06}

\begin{knowledgebox}{Lectura de la figura: es una compuerta del proceso de investigación, no una figura retórica de apertura}
Si a los usuarios les preocupa sobre todo la pérdida de la lengua en la educación, el sistema no puede optimizar solo la traducción de noticias; si la comunidad usa varios sistemas de escritura o variantes locales, unificarlos en un único estándar puede destruir directamente la utilidad. Por eso, las entrevistas sobre necesidades modifican el muestreo de datos, el sistema de etiquetas, los escenarios de evaluación y la release policy.
\end{knowledgebox}

\begin{importantbox}{Human-centered no equivale a hacer una prueba con usuarios al final}
La cadena de participación más completa es:
\[
\text{Entrevistas sobre necesidades}\rightarrow\text{Negociación del estándar lingüístico}\rightarrow\text{Traducción profesional}
\rightarrow\text{Evaluación por hablantes nativos}\rightarrow\text{Retroalimentación y corrección de errores}.
\]
Los usuarios no son inspectores de aceptación al final del pipeline, sino copartícipes en la definición del problema y en la producción de evidencia.
\end{importantbox}

\subsection{Muestra de entrevistas y límites del método}

El equipo entrevistó a 44 hablantes nativos de 36 lenguas y de varias regiones. Esta muestra no puede representar a todas las comunidades de lenguas de bajos recursos, pero basta para exponer problemas que una ruta puramente técnica no ve: la percepción del declive de la lengua, la presión de las lenguas dominantes en la educación, la expectativa de ser incluidos en los sistemas de traducción existentes y la necesidad real de algo «no perfecto, pero suficientemente útil».

\lecturefigure{slide-07-speaker-interviews.jpg}{Alcance de las entrevistas a hablantes nativos de lenguas de bajos recursos}{Video de la clase de Stanford 00:05:54}

\begin{knowledgebox}{Lectura de la figura: las cifras de la muestra indican el alcance exploratorio, no permiten inferencias sobre la población}
Los 44 entrevistados y las 36 lenguas muestran que el equipo recogió opiniones de manera deliberada entre lenguas distintas, pero a cada lengua le corresponden muy pocas personas, y la propia clase reconoce que los canales de reclutamiento tienen sesgos regionales y hacia grupos migrantes. El uso correcto es formular design hypotheses y seguir consultando en las etapas de traducción, evaluación y publicación, no presentar las proporciones de las entrevistas como conclusiones estadísticas globales.
\end{knowledgebox}

\teachervoice{Fan llama a este paso un social-sciences-type o sociology-type approach. No afirma que las entrevistas puedan sustituir a los experimentos, sino que, en problemas cercanos a las personas, solo después de aclarar el objetivo se sabe qué deben medir los experimentos posteriores.}

\subsection{Declive de la lengua, inclusión tecnológica y lo «suficientemente bueno»}

Los resultados de las entrevistas situaron la traducción automática en un contexto social más amplio: a algunos usuarios les preocupa el declive de la lengua local en la educación y en la vida pública; al mismo tiempo, ser incluidos en los sistemas de traducción dominantes genera fuertes expectativas. La clase señala además que los usuarios no necesariamente exigen que la traducción automática cumpla todos los estándares de la traducción humana profesional; en ciertos escenarios de comunicación cotidiana, una herramienta con defectos evidentes pero que va en la dirección correcta también puede aportar valor.

\lecturefigure{slide-08-interview-findings.jpg}{Preocupaciones, expectativas y juicios de utilidad revelados por las entrevistas}{Video de la clase de Stanford 00:07:03}

\begin{knowledgebox}{Lectura de la figura: los tres hallazgos actúan sobre tres niveles distintos del sistema}
El declive de la lengua afecta las prioridades del proyecto y el riesgo cultural; la expectativa de «ser incluidos en el sistema» afecta los objetivos de cobertura; que la traducción automática pueda ser ya suficiente en el uso real exige que la evaluación distinga escenarios como la publicación profesional, la educación, los viajes y la salud. Un puntaje promedio no puede sustituir estas distintas tolerancias al riesgo.
\end{knowledgebox}

\begin{warningbox}{«Suficientemente útil» no significa rebajar el estándar para las lenguas de bajos recursos}
Los umbrales por escenario permiten que tareas distintas tengan tolerancias al error distintas, pero eso no justifica dejar a las lenguas de bajos recursos para siempre con una calidad inferior. Sobre todo en información médica, de salud pública, legal y de crisis, el error de traducción de una palabra pequeña puede tener consecuencias graves; el ejemplo de toxicity más adelante lo mostrará en concreto.
\end{warningbox}

\teachervoice{En la sesión de preguntas y respuestas de la clase se enfatiza de nuevo que la consulta con hablantes nativos debe abarcar todo el proceso de desarrollo: desde las entrevistas iniciales, pasando por el equipo de traducción profesional, hasta la evaluación humana y la retroalimentación abierta de la comunidad. Human-centered es una propiedad de todo el ciclo de vida, no un workshop de una sola vez.}

\subsection{Resumen de la sección}

Antes de modelar, NLLB confirmó las necesidades reales y extendió la participación de la comunidad a los datos y la evaluación. Las entrevistas aportan evidencia de diseño e indicios de riesgo, no estadísticas representativas; lo «útil» debe definirse por escenario y no puede convertirse en excusa para eludir la responsabilidad sobre la calidad.
